% Canonical LaTeX source; imported and revised from the legacy Markdown.
\chapter{Studio di funzione e metodi numerici}
\label{chap:13}

\begin{obiettivocapitolo}{eseguire uno studio di funzione completo in ordine stabile}
\prioritaalta. La sequenza seguente riduce omissioni e contraddizioni tra grafico, limiti e derivate.
\end{obiettivocapitolo}
\begin{proceduraesame}{studio completo di funzione}
\begin{enumerate}
\item Dominio e punti esclusi.
\item Simmetrie, periodicità e intersezioni con gli assi.
\item Segno della funzione.
\item Limiti agli estremi del dominio e nei punti esclusi.
\item Asintoti verticali, orizzontali e obliqui.
\item Continuità e classificazione delle discontinuità.
\item Derivata prima, punti critici, tabella di monotonia ed estremi.
\item Derivata seconda, concavità/convessità e flessi con cambio di segno.
\item Valori notevoli e confronto con eventuali estremi assoluti.
\item Grafico coerente con tutte le informazioni precedenti.
\end{enumerate}
\end{proceduraesame}
\begin{sceltametodo}{zeri numerici}
Usa la bisezione quando disponi di continuità e cambio di segno e vuoi convergenza garantita. Usa Newton quando hai una buona stima iniziale, derivata non piccola e regolarità; controlla sempre il residuo e l'intervallo.
\end{sceltametodo}
\begin{errorefrequente}{punti critici incompleti}
Gli estremi possono trovarsi anche in punti non derivabili e agli estremi dell'intervallo, non soltanto dove $f'=0$.
\end{errorefrequente}
\begin{controllofinale}{grafico}
Ogni ramo deve rispettare dominio, segno, limiti, monotonia, concavità e asintoti. Un solo conflitto indica un errore precedente.
\end{controllofinale}

\section{Candidati agli estremi assoluti}\label{candidati-agli-estremi-assoluti}

Se \(f\in C([a,b])\) ed è derivabile salvo eventualmente in un insieme finito \(N\subset(a,b)\):

\[
\operatorname*{arg\,max}_{[a,b]} f,\quad
\operatorname*{arg\,min}_{[a,b]} f
\subseteq
\{a,b\}\cup\{x\in(a,b):f'(x)=0\}\cup N.
\]

\section{Asintoti}\label{asintoti}

Asintoto verticale. Basta un solo laterale infinito fra quelli ammessi dal dominio:

\[
\displaystyle
x=x_0
\quad\Longleftrightarrow\quad
\lim_{x\to x_0^-}f(x)=\pm\infty
\quad\text{oppure}\quad
\lim_{x\to x_0^+}f(x)=\pm\infty.
\]

Asintoto orizzontale nella direzione \(\sigma\in\{+,-\}\):

\[
\displaystyle
y=\ell
\quad\Longleftrightarrow\quad
\lim_{x\to\sigma\infty}f(x)=\ell\in\mathbb R.
\]

Asintoto obliquo nella direzione \(\sigma\in\{+,-\}\):

\[
\displaystyle
y=m_\sigma x+q_\sigma,
\]

\[
\displaystyle
m_\sigma=\lim_{x\to\sigma\infty}\frac{f(x)}{x},
\qquad
q_\sigma=\lim_{x\to\sigma\infty}\bigl(f(x)-m_\sigma x\bigr),
\]

con \(m_\sigma,q_\sigma\in\mathbb R\) e \(m_\sigma\ne0\).

\[
\displaystyle
\lim_{x\to\sigma\infty}
\bigl[f(x)-(m_\sigma x+q_\sigma)\bigr]=0.
\]

Approfondimenti: \href{https://ingegnerismo.it/matematica/asintoto/}{asintoto} e \href{https://ingegnerismo.it/matematica/studio-di-funzione/}{studio di funzione}.

\section{Punti stazionari ed estremi}\label{punti-stazionari-ed-estremi}

\[
\displaystyle
x_0\text{ stazionario}
\quad\Longleftrightarrow\quad
f'(x_0)=0.
\]

Test del segno di \(f'\), assumendo \(f\) continua in \(x_0\) e derivabile nei due intorni laterali puntati:

\[
\begin{aligned}
f':+\to- &\Longrightarrow x_0\text{ massimo locale},\\
f':-\to+ &\Longrightarrow x_0\text{ minimo locale},\\
f'>0\text{ sui due lati oppure }f'<0\text{ sui due lati}
&\Longrightarrow x_0\text{ non è un estremo}.
\end{aligned}
\]

Test della derivata seconda:

\[
\begin{aligned}
f'(x_0)=0,\ f''(x_0)>0
&\Longrightarrow x_0\text{ minimo locale},\\
f'(x_0)=0,\ f''(x_0)<0
&\Longrightarrow x_0\text{ massimo locale},\\
f'(x_0)=0,\ f''(x_0)=0
&\Longrightarrow \text{test inconcludente}.
\end{aligned}
\]

Test della prima derivata non nulla, nelle ipotesi di Taylor di ordine \(m\ge2\):

\[
\displaystyle
f'(x_0)=\cdots=f^{(m-1)}(x_0)=0,
\qquad
f^{(m)}(x_0)\ne0.
\]

\begin{longtable}[]{@{}ll@{}}
\toprule\noalign{}
Caso & Classificazione \\
\midrule\noalign{}
\endhead
\bottomrule\noalign{}
\endlastfoot
\(m\) pari, \(f^{(m)}(x_0)>0\) & minimo locale \\
\(m\) pari, \(f^{(m)}(x_0)<0\) & massimo locale \\
\(m\) dispari & non estremo; per \(m\ge3\), flesso stazionario \\
\end{longtable}

Approfondimenti: \href{https://ingegnerismo.it/matematica/punto-critico/}{punto critico}, \href{https://ingegnerismo.it/matematica/massimo/}{massimo} e \href{https://ingegnerismo.it/matematica/minimo/}{minimo}.

\section{Convessità, concavità e flessi}\label{convessituxe0-concavituxe0-e-flessi}

Per \(f:I\to\mathbb R\):

\[
\displaystyle
f\text{ convessa}
\quad\Longleftrightarrow\quad
f(\lambda x+(1-\lambda)y)
\le
\lambda f(x)+(1-\lambda)f(y)
\]

\[
\displaystyle
f\text{ concava}
\quad\Longleftrightarrow\quad
f(\lambda x+(1-\lambda)y)
\ge
\lambda f(x)+(1-\lambda)f(y),
\]

per ogni \(x,y\in I\) e \(\lambda\in[0,1]\).

Se \(f\) è due volte derivabile su \(I\):

\[
\begin{aligned}
f\text{ convessa su }I&\Longleftrightarrow f''\ge0\text{ su }I,\\
f\text{ concava su }I&\Longleftrightarrow f''\le0\text{ su }I.
\end{aligned}
\]

\[
\begin{aligned}
f''>0\text{ su }I&\Longrightarrow f\text{ strettamente convessa},\\
f''<0\text{ su }I&\Longrightarrow f\text{ strettamente concava}.
\end{aligned}
\]

Se \(f\) è continua in \(x_0\) e la concavità cambia attraversando il punto, si ha un flesso. In particolare, quando \(f''\) esiste nei due intorni laterali puntati:

\[
\displaystyle
f''\text{ cambia segno attraversando }x_0
\quad\Longrightarrow\quad
x_0\text{ è un flesso}.
\]

\[
\displaystyle
f''(x_0)=0
\quad\not\Longrightarrow\quad
x_0\text{ è un flesso}.
\]

Approfondimenti: \href{https://ingegnerismo.it/matematica/convessita/}{convessità}, \href{https://ingegnerismo.it/matematica/concavita/}{concavità} e \href{https://ingegnerismo.it/matematica/flesso/}{flesso}.

\section{Bisezione e metodo di Newton (estensione)}\label{bisezione-e-metodo-di-newton-estensione}

Nel seguito \(\alpha\) è uno zero di \(f\), \(x_n\) l'approssimazione al passo \(n\) ed \(e_n=x_n-\alpha\) l'errore vero.

Se \(f\in C([a,b])\) e \(f(a)f(b)<0\), il metodo di bisezione produce intervalli annidati contenenti uno zero. Dopo \(n\) bisezioni:

\[
b_n-a_n=\frac{b-a}{2^n},
\qquad
|x_n-\alpha|\le\frac{b-a}{2^{n+1}},
\]

dove \(x_n=(a_n+b_n)/2\) e \(\alpha\) è lo zero contenuto nell'intervallo.

Metodo di Newton:

\[
x_{n+1}=x_n-\frac{f(x_n)}{f'(x_n)},
\qquad f'(x_n)\ne0.
\]

Vicino a uno zero semplice \(\alpha\), sotto le usuali ipotesi di regolarità, la convergenza è quadraticamente locale:

\[
e_{n+1}
=
\frac{f''(\alpha)}{2f'(\alpha)}e_n^2
+o(e_n^2),
\qquad e_n=x_n-\alpha.
\]

Teoria: \href{https://ingegnerismo.it/matematica/metodo-di-bisezione/}{metodo di bisezione} e \href{https://ingegnerismo.it/matematica/metodo-di-newton/}{metodo di Newton}.

Esercizi svolti: \href{https://ingegnerismo.it/matematica/teorema-degli-zeri-esercizi/}{teorema degli zeri e bisezione}.
